%!TEX encoding = UTF8
%!TEX root = 0-notes.tex

\chapter{Structures algébriques}
\label{chap:structures}

% useful for U and U_n
\section{Groupe}

\dfn{groupe}{
	Un \emphindex{groupe} $\bigpar{G, \cdot}$ est un ensemble muni d'une loi satisfaisant les axiomes suivants.
	\begin{enumerate}[label={\bf A\arabic*.}]
		\item
		\textbf{La loi est interne.} Pour tout $g, h\in G$, $g\cdot h \in G$.
		\item
		\textbf{La loi est associative.} Pour tout $g, h, k \in G$, $(g\cdot h)\cdot k = g \cdot (h \cdot k)$.
		\item
		\textbf{Élément neutre.} Il existe un élément $e\in G$ neutre pour la loi : pour tout $g\in G$, $g \cdot e = e \cdot g = g$.
		\item
		\textbf{Inverse.} Tout élément est inversible : pour tout $g\in G$, il existe $g^{-1} \in G$ tel que $g \cdot g^{-1} = e = g^{-1} \cdot g$.
	\end{enumerate}
}{dfn:groupe}

\nt{
	Dans un groupe, l'inverse à gauche est égal à l'inverse à droite.
}

\notation{
	Il est fréquent d'omettre la loi dans un groupe, celle-ci étant unique. 
	Le produit $g \cdot h$ sera alors noté $gh$.
}

\dfn{commutativité}{
	Une loi $\cdot$ sur un ensemble $E$ est \emphindex{commutative} dès que
		\[ a \cdot b = b \cdot a \]
	pour tous les éléments $a, b \in E$.
}{dfn:commutativite}

\nomen{
	Une loi $\cdot$ commutative est généralement notée $+$, par analogie à l'addition qui est toujours commutative, alors que la multiplication ne l'est parfois pas (dans le cas des matrices, par exemple).
}

\exe{}{
	Montrer que $\bigpar{ Z ,+}$ est un groupe mais que $\bigpar{Z, \cdot}$ n'en est pas un.
}{exe:Z-group}{

}

\section{Anneau et corps}

\subsection{Anneau}

\dfn{anneau}{
	Un \emphindex{anneau} $\bigpar{ R, +, \cdot} $ est un ensemble ayant deux opérations satisfaisant les axiomes suivants.
	\begin{enumerate}
		\item
		$\bigpar{R, +}$ est un groupe abélien.
		\item
		La loi $\cdot$ est associative : pour tout $a, b, c\in R$, $a \cdot (b \cdot c) = (a \cdot b) \cdot c$.
		\item
		La loi $\cdot$ se distribue sur la loi $+$ : pour tout $a, b, c\in R$, 
		\begin{itemize}[label=$\bullet$]
			\item $(a+b) \cdot c = a \cdot c + b \cdot c$ ; et
			\item $a \cdot (b+c) = a \cdot b + a \cdot c$.
		\end{itemize}
	\end{enumerate}
}{dfn:anneau}

\nt{
	Attention : la structure $\bigpar{ R, \cdot}$ n'est pas forcément un monoïde : elle n'admet pas forcément d'élément neutre.
}

\nomen{
	Un anneau $R$ est dit
	\begin{enumerate}[leftmargin=80pt]
		\item[\emphindex{Unitaire}] s'il existe $1\in R$ élément neutre de la loi  $\cdot$.
		$1 \cdot u = u \cdot 1 = u$ pour tout $u\in R$.
		\item[\emphindex{Commutatif}] si la loi $\cdot$ est commutative.
		$a \cdot b = b \cdot a$ pour tout $a, b\in R$.
	\end{enumerate}
}

\ex{}{
	\begin{enumerate}
		\item
		$R = \Z$ est un anneau commutatif unitaire.
		\item
		$R = 2\Z$ est un anneau commutatif non unitaire.
		\item
		$R = \{ 0 \}$ est un anneau commutatif unitaire en identifiant $1=0$.
		\item
		$R = \Z[i] = \bigset{ a + bi \tq a, b\in\Z}$, les \emphindex{entiers de Gauss}, est un anneau commutatif unitaire.
		\item
		$R = \Z^{n \times n}$ et $\R^{n\times n}$ sont des anneaux unitaires non commutatifs.
		\item
		$R = C^{0}(\R) = \bigset{ f : \R \rightarrow \R \tq \text{$f$ continue} }$ est un anneau commutatif unitaire.
	\end{enumerate}
}{ex:anneaux}

\nt{
	L'élément 0 étant défini pour l'addition, il n'est pas clair \apriori~ que multiplier par 0 donne 0.
	L'élément 1 étant défini pour la multiplication, il n'est pas clair que multiplier par son inverse additif aboutit à l'inverse additif du facteur.
	
	La proposition suivante permet de déduire des axiomes de l'anneau unitaire les propriétés intuitives de 0 et de 1.
}

\prop{propriétés de $0$ et $1$}{
	Soit $R$ un anneau unitaire. Alors, pour tout $u\in R$,
	\begin{enumerate}
		\item $u \cdot 0 = 0 \cdot u = 0$ ; et
		\item $u \cdot (-1) = (-1) \cdot u = -u$.
	\end{enumerate}
}{prop:prop-0-1}

\exe{}{
	Montrer que $u \cdot 0 = 0$ pour tout $u$ élément d'un anneau unitaire $R$.
}{exe:utimes0}{
	Nous avons la suite d'égalités suivantes.
		\begin{align*}
			u \cdot 0 &= u \cdot (0+0)  && \text{car $\bigpar{R, +}$ est un groupe d'élément neutre $0$} \\
				&= u \cdot 0 + u \cdot 0 && \text{par associativité de $ \cdot $ sur $+$} \\
			0 &= u \cdot 0 && \text{en soustrayant $u \cdot 0$}
		\end{align*}
}

\exe{}{
	Montrer que $u \cdot (-1) = -u$ pour tout $u$ élément d'un anneau unitaire $R$.
}{exe:utimes-1}{
	Nous avons la suite d'égalités suivantes.
		\begin{align*}
			u \cdot (-1) + u &= u \cdot (-1+1) && \text{par associativité de $ \cdot $ sur $+$} \\
				&= u \cdot 0 && \\
				&= 0 &&
		\end{align*}
}

\dfn{unité, ensemble des unités}{
	Un élément $u\in R$ est dit \emphindex{unité} de l'anneau $R$ unitaire s'il existe $v\in R$ tel que
		\[ u \cdot v = 1 = v \cdot u. \]
	L'\emphindex{ensemble des unités} est donné par
		\[ R^\times = \bigset{ u \in R \tq \text{$u$ est unité} }. \]
}{dfn:unite}

\nomen{
	Une unité est aussi appelé un élément \emphindex{inversible}.
}

\exe{}{
	Montrer que $\bigpar{ R^\times,  \cdot  }$ est un groupe.
}{exe:units-group}{
	todo
}

\subsection{Corps}

\dfn{corps}{
	Soit $R$ un anneau commutatif unitaire dans lequel on suppose $1\neq0$.
	$R$ est un \emphindex{corps} si $R^\times = R-\{0\}$.
	Chaque élément du corps, hormis 0, est inversible.
}{dfn:corps}

\notation{
	Nous noterons $\K$ les corps dans la suite.
}

\exe{}{
	Montrer que si $1=0$ dans un anneau $R$ commutatif unitaire, alors $R = \bigset{ 0 }$.
	
	Montrer que si $0$ est inversible, alors $1=0$.
}{exe:1-is-0}{
	Tout élément $u$ de $R$ vérifie $u = u \cdot 1 = u \cdot 0 = 0$.
	
	Si $0$ est inversible, alors $0 = 0 \cdot 0^{-1} = 1$.
}

\exe{}{
	Soit $p\in\N$ premier.
	Montrer que $\F_p = \bigpar{ \Z/p\Z, +, \cdot }$ est un corps.
}{exe:Fp-field}{
	todo
}

\exe{difficulty=1}{
	Soit $R$ un anneau commutatif unitaire et $\circ$ la loi définie par
		\[ a \circ b = a + b - ab. \]
	\begin{enumerate}
		\item 
		Montrer que $\circ$ est associative d'élément neutre 0.
		\item
		Montrer que $\bigpar{R-\{1\}, \circ}$ est un groupe si et seulement si $R$ est un corps.
	\end{enumerate}
}{exe:circle-op}{

}

\dfn{corps gauche}{
	Soit $R$ un anneau unitaire non commutatif dans lequel on suppose $1\neq0$.
	$R$ est un \emphindex{corps gauche} si $R^\times = R-\{0\}$.
	Chaque élément du corps, hormis 0, est inversible.
}{dfn:corps-gauche}

\ex{}{
	L'ensemble $\H$ des \emphindex{quaternions} défini par
		\[ \H = \bigset{ a+bi+cj+dk \tq a, b, c, d \in \R, i^2 = j^2 = k^2 = -1, \et ij = k, jk = i, ki = j} \]
	est un corps gauche.
	
	En effet, $(a+bi+cj+dk)^{-1} = \frac{a-bi-cj-dk}{a^2 + b^2 + c^2 + d^2}$ si $(a,b,c,d)\neq(0,0,0,0)$
	et $ij = (jk)(ki) = jk^2 i = -ji$.
	
	Il s'avère que les quaternions traduisent des rotations dans $\R^3$ et s'identifient à un ensemble de matrices $3\times3$.
	Leur non commutativité n'est donc pas surprenante.
}{ex:quaternions}

\subsection{Anneau intègre}

\nomen{
	Un élément non nul $u\in R-\{0\}$ d'un anneau $R$ est un \emphindex{diviseur de zéro} s'il existe $v\in R-\{0\}$ non nul également tel que $u \cdot v = 0$.
}

\exe{}{
	Soit $n\in\N^*$ et $R = \Z/n\Z$ l'anneau des entiers modulo $n$.
	Donner l'ensemble des diviseurs de zéro et l'ensemble des éléments inversibles de $R$.
}{exe:ZnZ-div-zero}{
	Nous allons montrer que $R$ se partionne comme
		\[ R = \{ 0 \} \sqcup R^\times \sqcup \bigset{ \text{ diviseurs de zéro } }. \]
	D'abord, un élément $x\in R-\{0\}$ est diviseur de zéro si et seulement si $d = \pgcd(x,n) \neq 1$ dans $\Z$.
	En effet, si $x$ divise zéro, alors il existe un $y \in R-\{0\}$ tel que $xy = 0$.
	Ainsi, $n$ divise le produit $xy$ mais $n$ ne divise pas $y$.
	Ainsi $d\neq1$, sinon le lemme d'Euclide impliqerait $n|y$.
	
	Réciproquement, si $d \neq 1$, alors $x \times \frac{n}d = \frac{x}d n = 0$ dans $R$, et $\frac{n}d \neq 0$.
	
	D'abord, un élément $x\in R-\{0\}$ est inversible si et seulement si $d = \pgcd(x,n) = 1$ dans $\Z$.
	Si $x$ est inversible, il existe $y\in R$ tel que $xy = 1$.
	Nous obtenons une identité de Bézout : $xy + un = 1$ pour un certain $u\in\Z$.
	Comme le PGCD $d$ divise $x$ et $n$ dans $\Z$, il divise donc aussi 1, et $d=1$.
	
	Réciproquement, si $d=1$. alors l'identité de Bézout $xu + nv = 1$ avec $u, v \in \Z$ implique que $xu = 1$ dans $R$ et que $x$ est unité.
}

\exe{}{
	Considérons $R = \bigset{ f : \R \rightarrow \R }$ l'ensemble des applications de  $\R$ dans $\R$.
	\begin{enumerate}
		\item Donner les éléments neutres pour l'addition et la multiplication dans l'anneau $\bigpar{R, +, \cdot}$.
		\item Décrire les diviseurs de zéro de $R$.
		\item Montrer que chaque élément non nul est soit un diviseur de zéro, soit inversible.
	\end{enumerate}
}{exe:ring-of-maps}{
	todo
}

\prop{}{
	Soit $R$ un anneau unitaire et $u\in R$ une unité.
	Alors $u$ ne divise pas zéro.
}{prop:unit-dont-divide-zero}

\pf{}{
	Si $u$ divise zéro, alors il existe $v\neq0$ tel que $u\cdot v = 0$.
	Comme $u$ est unité, nous avons $v = (u^{-1} u) v = u^{-1} (uv) = u^{-1} 0 = 0$.
	Ceci est contradictoire.
}

\nomen{
	Un anneau $R$ commutatif unitaire est $\emphindex{intègre}$ dès qu'il n'admet aucun diviseur de zéro.
}

\exe{}{
	Montrer que la réciproque de la proposition \ref{prop:unit-dont-divide-zero} est fausse.
	Au regard de l'exercice \ref{exe:ZnZ-div-zero}, considérer un anneau de caractéristique nulle.
}{exe:notunit-nordivzero}{
	L'élément 2 de l'anneau $\Z$ n'est pas une unité ni un diviseur de zéro.
}

\cor{}{
	Tout corps $\K$ est intègre : si $u \cdot v = 0$ dans $\K$, alors $u = 0$ ou $v = 0$.
}{cor:corps-integre}

\subsection{Idéal}

\dfn{idéal}{
	Soit $R$ un anneau et $I\subset R$ un ensemble.
	\begin{enumerate}
		\item
		$I$ est un \emphindex{idéal à gauche} si
		\begin{enumerate}
			\item
			$\bigpar{I, +}$ est un groupe abélien ; et
			\item
			$r \cdot  I \subset I$ pour tout $r\in R$.
		\end{enumerate}
		\item
		$I$ est un \emphindex{idéal à droite} si
		\begin{enumerate}
			\item
			$\bigpar{I, +}$ est un groupe abélien ; et
			\item
			$I \cdot r \subset I$ pour tout $r\in R$.
		\end{enumerate}
		\item
		$I$ est un \emphindex{idéal} ou \emphindex{idéal bilatère} s'il est un idéal à gauche et à droite.
	\end{enumerate}
}{dfn:ideal}

\nt{
	Attention : lorsque non précisé « à gauche » ou « à droite », un idéal est donc un idéal bilatère.
}

\nt{
	Un idéal est un sous-anneau non unitaire qui est stable par multiplication extérieure.
}

\notation{
	L'expression « $I$ est un idéal de $R$ » est notée $I \lhd R$.
}

\exe{}{
	Soit $I, J \subseteq R$ deux idéaux d'un anneau $R$.
	Montrer que les ensembles suivants sont également des idéaux de $R$ et que chaque ensemble est inclus dans le suivant.
	\begin{enumerate}
		\item l'idéal produit $I\cdot J = \bigset{ a_1 b_1 + \dots + a_n b_n \tq a_i \in I, b_i \in J, n\in\N^* }$ ;
		\item l'idéal intersection $I\cap J$ ;
		\item l'idéal somme $I+J = \bigset{ a + b \tq a\in I, b\in J}$.
	\end{enumerate}
}{exe:somme-produit-inter-ideaux}{
	todo
}

\prop{}{
	Soit $I\lhd R$ un idéal d'un anneau. Alors
	\begin{enumerate}[label=(\roman*)]
		\item
		$I=R$ si et seulement si $I$ contient une unité.
		\item
		si $R$ est commutatif, $R$ est un corps si et seulement si ses seuls idéaux sont $\{ 0 \}$ et $R$.
	\end{enumerate}
	L'étude des idéaux n'est donc pas très intéressante dans un corps !
}{prop:ideaux-corps}


\dfn{idéal généré par un ensemble}{
	Soit $R$ un anneau et $A \subseteq R$ un sous-ensemble.
	L'\emphindex{idéal engendré par $A$}, noté $\japb{A}$, est le plus petit idéal de $R$ contenant $A$.
}{dfn:ideal-gen}

\prop{}{
	Soit $R$ un anneau et $A \subseteq R$ un sous-ensemble.
	Alors
		\[ \japb{A} = \Bigset{ \sum_{i=1}^n r_i a_i r_i' \tq r_i, r_i' \in R, a_i \in A, n\in\N^*}. \]
	Si $R$ est commutatif, alors
		\[ \japb{A} = \Bigset{ \sum_{i=1}^n r_i a_i \tq r_i \in R, a_i \in A, n\in\N^*}. \]
}{prop:ideal-gen}

\nomen{
	Un idéal est de \emphindex{type fini} s'il est généré par un ensemble fini.
	Un idéal généré par un unique élément est \emphindex{principal} et noté $I = \japb{a} = (a)$.
}

\exe{}{
	Soit $a, b\in\Z$ deux entiers relatifs.
	Montrer l'équivalence suivante.
		\begin{align*}
			(a) \subseteq (b) \iff b | a
		\end{align*}
}{exe:incl-div}{
	todo
}

\exe{difficulty=1}{
	Soit $I \lhd \Z$ un idéal et $d \in I$ élément de l'idéal de valeur absolue minimale.
	Montrer que $d$ divise tous les éléments de $I$ et donc que $I = \japb{d}$.
}{exe:Z-principal}{
	todo
}

\section{Algèbre sur un corps}

% for  poly(x)

\section*{Exercices}
\addcontentsline{toc}{section}{Exercices}

\shipoutExercise

\section*{Solutions}
\addcontentsline{toc}{section}{Solutions}

\shipoutAnswer



